% vkSplat: A Compute-First Vulkan Path for Robotics Synthetic Data
% Position / vision paper draft (v0.1)
% Target venue: HotOS-style position track, HPG, MLSys, or arXiv preprint.

\documentclass[10pt,journal]{IEEEtran}

\usepackage{graphicx}
\usepackage{amsmath,amssymb}
\usepackage{booktabs}
\usepackage{hyperref}
\usepackage{cite}
\usepackage{xcolor}
\usepackage{url}

\newcommand{\todo}[1]{\textcolor{red}{[TODO: #1]}}
\newcommand{\note}[1]{\textcolor{blue}{[NOTE: #1]}}

\title{vkSplat: A Compute-First Vulkan Path for Robotics\\Synthetic Data Generation on AI Accelerators}

\author{%
    Chien-Ping (CP) Lu\\%
    \textit{Independent}\\%
    \texttt{chienpng.lu@gmail.com}%
}

\begin{document}
\maketitle

\begin{abstract}
Modern robotics increasingly depends on synthetic data generated inside
photorealistic digital twins. The dominant pipelines --- NVIDIA Isaac~Sim
\slash Omniverse, Cosmos, BlenderProc, Kubric --- bind the rendering layer to
graphics-optimized silicon: GPUs with rasterizers, ROPs, and ray-tracing
cores. We argue that this binding is incidental rather than fundamental.
The trend is already visible inside graphics itself: DLSS and its peers have
shifted the center of gravity of real-time rendering from fixed-function
graphics units to \emph{tensor compute}, with the rasterizer and RT cores
producing only a low-resolution, temporally and spatially sparse seed that
neural networks then upscale, denoise, anti-alias, and
interpolate~\cite{nvidia-dlss,liu2020nsrr,nvidia-dlss3,schied2017svgf,bitterli2020restir}.
Synthetic data generation (SDG) is a \emph{throughput-bound,
latency-tolerant} workload, and its rendering quality requirements --- while
non-trivial --- are met by published compute-only techniques (software
rasterization, GPU ray traversal, 3D Gaussian splatting) that do not require
fixed-function graphics hardware. We propose \textbf{vkSplat}: a Vulkan
implementation in which the API surface is preserved as a portable contract,
but the backend is a software renderer expressed in CUDA-class compute. The
near-term motivation is access: data-center AI accelerators (H100, B200, AI
chips such as TPU and Trainium) are over an order of magnitude more available,
in aggregate FLOPS and procurement velocity, than RTX-class workstation
silicon. The longer-term motivation is portability: decoupling the API from
the silicon turns digital-twin pipelines into a software problem with a clean
backend interface. We present the system design, situate it relative to
existing software-Vulkan implementations (SwiftShader, Mesa Lavapipe), justify
3D Gaussian Splatting as the v1 renderer, and lay out an evaluation plan for
throughput, energy, and downstream sim-to-real task accuracy. \todo{numbers}
\end{abstract}

\begin{IEEEkeywords}
Synthetic data generation, Vulkan, software rendering, 3D Gaussian splatting,
digital twins, robotics, AI accelerators.
\end{IEEEkeywords}

% =====================================================================
\section{Introduction}
% =====================================================================

Two trends are colliding. First, robotics has become a data-hungry
discipline: end-to-end visuomotor policies, foundation models for
manipulation, and autonomy stacks all require enormous quantities of
labeled visual data, and an increasing fraction of that data is
\emph{synthetic} --- rendered inside high-fidelity simulators
\cite{tobin2017domain,greff2022kubric,denninger2019blenderproc,deitke2022procthor}.
Second, the silicon supply chain has bifurcated: AI accelerators
(NVIDIA~H100\slash B200, Google~TPU, AWS Trainium, and a long tail of
emerging chips) are scaling FLOPS and memory bandwidth at a pace that
outstrips the workstation-class GPUs traditionally used for graphics, and
they dominate procurement budgets at every hyperscaler.

The collision is awkward. Today's leading SDG stacks --- NVIDIA
Omniverse\slash Isaac~Sim\slash Replicator and Cosmos
\cite{nvidia2024cosmos}, Unreal-based platforms, Blender-based pipelines
\cite{denninger2019blenderproc,greff2022kubric} --- are built on graphics
APIs (Vulkan, D3D12, Metal, OpenGL) whose performant implementations
require dedicated graphics silicon: rasterizer pipelines, ROPs, RT cores.
The data-center accelerators that companies actually have are largely
graphics-blind: an H100 has no RT cores, and TPUs and Trainium have no
graphics pipeline at all. Generating training data on the chips one
already owns therefore requires either (a) provisioning a separate fleet
of graphics GPUs or (b) abandoning the graphics-API ecosystem in favor of
custom rendering harnesses, which fragments the tooling.

This paper takes the position that this fragmentation is unnecessary.
Vulkan is a clean, modern, vendor-neutral API. There is no fundamental
obstacle to implementing it as a \emph{software renderer expressed in
compute kernels}, lowering rasterization, ray tracing, and shadow mapping
to general-purpose accelerator compute. Software Vulkan implementations
already exist (SwiftShader~\cite{swiftshader}, Mesa
Lavapipe\slash llvmpipe~\cite{mesa-lavapipe}), but they target CPUs for
correctness and conformance, not GPUs for throughput; they do not exploit
the parallelism of AI accelerators, and they have not been engineered for
SDG workloads. We propose to fill that gap.

The argument has, in fact, already been made by the
graphics-hardware industry's own roadmap. In modern real-time pipelines
exemplified by NVIDIA DLSS~\cite{nvidia-dlss,liu2020nsrr,nvidia-dlss3},
the rasterizer and RT cores no longer produce the final image: they
produce a \emph{low-resolution, temporally jittered, spatially sparse,
under-sampled path-traced seed}. The rest of the pipeline ---
super-resolution, anti-aliasing, denoising, frame
interpolation~\cite{schied2017svgf,chaitanya2017interactive,nvidia-nrd},
and importance-resampled direct illumination~\cite{bitterli2020restir} ---
runs on tensor units. Each generation of DLSS shrinks the seed and
expands the neural reconstruction. \emph{The center of gravity of the
real-time rendering pipeline has already moved off fixed-function
graphics hardware and onto tensor compute.} What we propose is the
logical next step: if the AI portion is doing the heavy lifting, replace
the fixed-function seed generator with a software-rendered seed produced
on the same tensor-rich silicon, and eliminate the dependence on
graphics-specific units altogether. For SDG, where real-time interaction
is unnecessary, this trade is straightforward.

\textbf{Contributions.} This is a position paper introducing the
\textbf{vkSplat} project. We make four claims:

\begin{enumerate}
    \item \textbf{Workload.} Robotics SDG is throughput-bound and
    latency-tolerant. A renderer that produces images in tens to hundreds
    of milliseconds is acceptable if it sustains high aggregate throughput
    and high quality. This relaxes the constraint that historically
    forced the use of fixed-function graphics hardware
    (Section~\ref{sec:workload}).

    \item \textbf{API decoupling.} The Vulkan API can be preserved as the
    application contract while the implementation is replaced by a
    compute-only software renderer. We outline the architecture of such
    an implementation (Section~\ref{sec:design}).

    \item \textbf{3DGS-first renderer.} 3D Gaussian
    Splatting~\cite{kerbl20233dgs} is the natural v1 renderer because it
    is already compute-native: no rasterizer, no ROPs, no RT cores
    required. We argue this lets the project ship a usable system before
    tackling the harder problem of lowering classical raster and ray
    tracing (Section~\ref{sec:v1}).

    \item \textbf{Hardware-agnostic backend.} The compute IR underlying
    the renderer can be retargeted from CUDA to SYCL, Triton, or
    AI-accelerator native APIs, broadening deployment to any chip with
    sufficient compute and memory bandwidth (Section~\ref{sec:backends}).
\end{enumerate}

\textbf{Non-goals.} vkSplat is not aimed at real-time gaming, VR, or
display-bound interactive workloads. We do not claim to outperform RTX
hardware on its own ground. We claim only that, for the SDG workload,
running graphics on AI compute is viable, economically attractive, and
strategically preferable.

% =====================================================================
\section{Background}\label{sec:bg}
% =====================================================================

\subsection{Synthetic data generation for robotics}

The use of simulated environments to train physical agents predates deep
learning, but its modern form was crystallized by domain
randomization~\cite{tobin2017domain} and by simulation platforms such as
Habitat~\cite{savva2019habitat}, ThreeDWorld~\cite{gan2020threedworld},
ProcTHOR~\cite{deitke2022procthor}, and NVIDIA Isaac~Sim. Recent SDG
toolkits like BlenderProc~\cite{denninger2019blenderproc} and
Kubric~\cite{greff2022kubric} provide higher-level scripting for
photorealistic image and video generation with controlled randomization
of geometry, materials, and lighting. NVIDIA Cosmos~\cite{nvidia2024cosmos}
extends this trajectory by combining classical rendering with generative
priors. Across these systems, the rendering layer is the single largest
silicon dependency.

\subsection{Software Vulkan implementations}

Two production software-Vulkan implementations exist.
\textbf{SwiftShader}~\cite{swiftshader}, originally a software renderer
project at TransGaming and now maintained by Google, provides a portable,
LLVM-JIT-based Vulkan~1.3 implementation that runs on CPUs (it is shipped
in Chromium for fallback rendering). \textbf{Mesa Lavapipe}~\cite{mesa-lavapipe}
is the Vulkan front-end on top of Mesa's \texttt{llvmpipe} CPU rasterizer.
Both are correctness-first and CPU-targeted; their performance on GPUs is
either non-existent (Lavapipe is CPU-only) or limited (SwiftShader has
experimental Reactor backends but is not optimized for AI-accelerator
throughput). Neither targets the SDG workload class.

\subsection{Compute-only graphics on GPUs}

There is a long line of research in software rendering on
general-purpose GPU compute. Laine and
Karras~\cite{laine2011software} showed that a fully software rasterizer
on CUDA could approach hardware rasterizer throughput for many workloads,
particularly when rasterization was not the bottleneck.
For ray tracing, Aila and Laine~\cite{aila2009understanding} established
the architectural patterns that would later become OptiX~\cite{parker2010optix}
and Embree~\cite{wald2014embree}. Path tracers including
Mitsuba~3~\cite{jakob2022mitsuba3} (built on Dr.Jit~\cite{jakob2022drjit})
and Falcor~\cite{benty2020falcor} demonstrate that high-quality,
differentiable rendering can be expressed entirely in compute.
Meanwhile, 3D Gaussian Splatting~\cite{kerbl20233dgs} and its
descendants~\cite{yu2024mip,huang20242dgs} have established a new class
of renderer that is compute-native by construction.

The intellectual point is simple: every major step of a modern graphics
pipeline has a published, competitive software-on-compute implementation.
\textbf{What does not exist is a unified Vulkan front-end that exposes
these implementations through a standard API to existing graphics
applications.} That is the gap vkSplat targets.

% =====================================================================
\section{Why SDG is the right workload}\label{sec:workload}
% =====================================================================

\subsection{Throughput-bound, latency-tolerant}

A real-time renderer must finish a frame in $\sim$16~ms or it fails.
A SDG pipeline must produce a target number of labeled images per dollar
or per joule; per-image latency is bounded only by the human iteration
loop and downstream training pipeline. If a software renderer takes
50~ms per image but uses one-tenth the silicon cost of an RTX-class
solution, it wins on $\$\slash\textrm{image}$ even though it loses on
frame time.

\todo{Add a table: typical SDG image budgets vs.\ real-time budgets;
illustrative throughput numbers.}

\subsection{Quality requirements}

SDG renderers do not need to convince a human that the image is real;
they need to teach a model that the image is informative. Modern
research~\cite{tobin2017domain,nvidia2024cosmos} suggests that
\emph{diversity}, \emph{labeling fidelity}, and \emph{controlled
randomization} matter more than the marginal photorealism that justifies
RT-core hardware. A 3DGS rendering of a real scene captured with a
phone, augmented with diverse synthetic objects, can be more useful for
training than a mathematically prettier path-traced render of an
unrealistic asset library.

\subsection{Hardware availability and cost}

\todo{Add a quantitative argument: aggregate H100\slash B200\slash TPU
FLOPS available globally vs.\ aggregate L40\slash RTX-6000-Ada FLOPS;
spot pricing per FLOP\slash s\slash dollar; energy per image budget.}

The qualitative point is that any team that wants to do SDG today is
either competing with rendering studios for L40s and RTX 6000 Adas, or
running their renderer on consumer hardware that is unstable in
data-center deployments. AI compute is, in contrast, the silicon those
teams already have purchased for training.

\subsection{The DLSS precedent: tensor compute is already where the
rendering value lives}\label{sec:dlss}

The strongest evidence that graphics workloads can move off
graphics-specific silicon is that they already have, partially, inside
NVIDIA's own products. We summarize the trajectory:

\begin{itemize}
    \item \textbf{DLSS~1 (2018).} Per-game neural super-resolution
    networks reconstruct high-resolution frames from sub-native input.
    The renderer's job is reduced from ``produce a $4\textrm{K}$
    image'' to ``produce a $1080\textrm{p}$ image.''

    \item \textbf{DLSS~2 (2020)~\cite{liu2020nsrr,edelsten2019dlss}.}
    Generic temporal super-resolution. The rasterizer's frame budget
    drops further and motion vectors become a first-class input. AI is
    now responsible for both spatial and temporal reconstruction.

    \item \textbf{DLSS~3 / Frame
    Generation~(2022)~\cite{nvidia-dlss3}.} Optical-flow-conditioned
    frame interpolation generates entire intermediate frames on tensor
    cores --- frames the rasterizer never drew at all. The
    fixed-function pipeline now contributes only a fraction of the
    pixels the user sees.

    \item \textbf{Ray-traced global illumination workflows} have
    converged on the same pattern. Single-sample-per-pixel path
    tracing, denoised by spatiotemporal
    filters~\cite{schied2017svgf,chaitanya2017interactive,nvidia-nrd}
    and importance-resampled by ReSTIR~\cite{bitterli2020restir}, has
    replaced classical many-sample Monte Carlo. RT cores produce a few
    sparse samples; tensor and shader cores reconstruct the rest.
\end{itemize}

The implication is direct. The fixed-function units already do not
produce the final image; they produce a sparse, low-resolution, noisy
\emph{seed} that AI then upscales, denoises, and interpolates. As the
seed shrinks and the AI grows, the share of total rendering work done
on tensor cores rises monotonically, and the marginal value of
fixed-function silicon falls.

Two questions then arise. \textbf{(1)} If the seed is small and noisy
anyway, can a software renderer on compute produce it at acceptable
quality? Published work on software rasterization~\cite{laine2011software},
GPU ray traversal~\cite{aila2009understanding}, and compute-native 3DGS
~\cite{kerbl20233dgs} suggests yes. \textbf{(2)} Once the seed is on
compute, does any reason remain to maintain a separate fixed-function
data path? For SDG, where there is no real-time constraint and where
the final consumer of the image is a learned model rather than a human
eye, the answer is no.

vkSplat is the experiment that pushes this thesis to its conclusion.
We do not argue that fixed-function graphics will disappear from
real-time gaming. We argue that for the throughput-bound,
latency-tolerant SDG workload, the DLSS-era trajectory has already made
fixed-function graphics optional, and the implementation should be
redrawn accordingly.

% =====================================================================
\section{System design}\label{sec:design}
% =====================================================================

vkSplat is structured as three layers
(Fig.~\ref{fig:arch})\todo{draw figure}:

\begin{enumerate}
    \item \textbf{Vulkan front-end.} A Vulkan loadable ICD that
    implements a \emph{conformance-targeted subset} of Vulkan~1.3 plus
    selected ray-tracing and external-memory extensions. Applications
    link the standard Vulkan loader and need no modification.

    \item \textbf{Intermediate representation (IR).} SPIR-V shaders and
    Vulkan command buffers are lowered into a backend-agnostic compute
    IR. We propose to base this IR on existing infrastructure ---
    LLVM\slash MLIR for shader code, a record-and-replay command-buffer
    IR for pipeline state --- so that the system rides on existing
    compiler ecosystems rather than inventing one.

    \item \textbf{Compute backend.} The IR is lowered to one of several
    backends: CUDA (reference), Triton or SYCL for portability, and
    eventually AI-accelerator-native paths (XLA HLO for TPU, NeuronCore
    for Trainium, etc.). The backend implements the rendering kernels:
    rasterization, ray traversal, shadow mapping, and
    Gaussian-splat blending.
\end{enumerate}

\subsection{Vulkan front-end}

We target a deliberate \emph{subset}: graphics pipelines, compute
pipelines, command buffers, descriptor sets, swapchain (optional, for
debug viewer), timeline semaphores, external memory. We initially omit
multi-draw-indirect at high feature parity, mesh shaders, sparse
binding, and presentation-engine optimization. SDG does not need them.

\subsection{Lowering rasterization}

We lower the classical raster pipeline to a tile-based software
rasterizer in the spirit of Laine and
Karras~\cite{laine2011software}, with each output tile assigned to a
threadblock\slash workgroup that performs binning, depth test, and
fragment shading.

\subsection{Lowering ray tracing and shadow mapping}

Ray-tracing pipelines (\texttt{VK\_KHR\_ray\_tracing\_pipeline}) lower to
software BVH traversal kernels following the
Aila--Laine~\cite{aila2009understanding} pattern, optionally calling
into Embree-like tree construction routines on the host. Shadow maps
become depth-only raster passes followed by a comparison kernel; for
high-quality shadows we expect ray-traced shadows to be cheaper than
classical shadow-map filtering at SDG quality targets.

\subsection{The 3DGS fast path}\label{sec:v1}

We expose a non-standard Vulkan extension
(\texttt{VK\_VKSPLAT\_gaussian\_splatting}) that lets applications submit
Gaussian primitives directly. Internally, this dispatches to a
3DGS rasterizer akin to \cite{kerbl20233dgs}'s but extended for offline
SDG (deterministic ordering, deeper sample budgets per pixel, optional
denoising). The rationale is that 3DGS does not need any of the
classical pipeline machinery, and a system that ships a working 3DGS
backend on day one is immediately useful for capturing real environments
and rendering them as training scenes.

% =====================================================================
\section{Backends}\label{sec:backends}
% =====================================================================

\textbf{CUDA.} The reference implementation, targeting NVIDIA H100 and
B200. CUDA is mature, has a rich kernel ecosystem (cuBLAS, CUB, Thrust,
OptiX-as-library), and provides external-memory and external-semaphore
interop with Vulkan~\cite{nvidia2020vkcuda}.

\textbf{Triton\slash SYCL.} A second backend in Triton or SYCL would
demonstrate portability and allow deployment on AMD MI300, Intel Ponte
Vecchio, and any future accelerator with an LLVM-class compiler.

\textbf{AI-accelerator-native.} TPU (XLA HLO), Trainium (NeuronCore SDK),
Cerebras, and Groq each have native compilation paths. Lowering rendering
kernels to these stacks is non-trivial but, for \emph{compute-dominant}
kernels (3DGS rasterization, BVH traversal), feasible. The rendering
problem becomes structurally close to a sparse-attention problem, which
these chips already optimize for.

% =====================================================================
\section{Evaluation plan}
% =====================================================================

We will evaluate vkSplat along three axes.

\textbf{Throughput.} Images per second on Mip-NeRF~360, Tanks~\&
Temples, and Deep Blending scenes. \todo{add SDG-specific scene
benchmarks: ProcTHOR rooms, Isaac Lab industrial scenes.}

\textbf{Cost \& energy.} Dollars-per-image and joules-per-image on a
range of accelerators (H100, B200, A100, L40, MI300, TPU~v5e). The L40
serves as a graphics-hardware baseline; the rest are pure compute.

\textbf{Downstream task accuracy.} Train a small object-detection or
manipulation policy on data generated by vkSplat vs.\ data generated by
Isaac~Sim or BlenderProc, controlling for total training compute. The
hypothesis is that vkSplat data is at least competitive on per-dollar
performance.

\todo{Specify benchmark scenes, target metrics, and ablation matrix in
detail.}

% =====================================================================
\section{Discussion}
% =====================================================================

\textbf{When does this lose?} If the workload is dominated by
fixed-function operations (heavy alpha-tested foliage, complex
post-processing chains, high-fidelity displacement), software renderers
on compute will pay a large constant factor. Real-time interactive use
will stay on RTX hardware for the foreseeable future.

\textbf{Why not a new API?} A new graphics API for AI compute has been
proposed informally many times. The reason to keep Vulkan is ecosystem:
existing engines (Unreal's Vulkan RHI, Unity, Godot, Filament), tooling
(RenderDoc, NSight, ImGui), and middleware all speak Vulkan. Reusing
that surface means a vkSplat backend benefits the ecosystem
asymmetrically.

\textbf{Strategic implications.} If a viable Vulkan-on-AI-compute path
exists, the SDG market is no longer constrained by the supply of
graphics silicon. This is structurally similar to what NVIDIA Omniverse
does for digital twins, but with the asset-distribution and
hardware-binding stories inverted: vkSplat is open at the API, neutral
at the hardware, and assumes a software-defined data center as its
deployment target. It is also consistent with the direction the
graphics-hardware industry has been moving for half a decade
(Section~\ref{sec:dlss}): each successive generation of DLSS-class
neural reconstruction has shifted more of the rendering pipeline onto
tensor cores. vkSplat is the architectural endpoint of that
trajectory --- a renderer in which \emph{all} of the work, seed and
reconstruction alike, runs on tensor-rich compute, and graphics
hardware is no longer in the loop.

\textbf{Limitations.} v1 will not be Vulkan-conformant. It will be
correct on a defined subset, on defined drivers, for defined SDG
workloads. Conformance is a long-term goal but not a pre-condition for
usefulness.

% =====================================================================
\section{Related work, expanded}
% =====================================================================

\todo{Expand: tile-based GPU rasterization (Liu et al.), differentiable
rendering (Loper~\&~Black~2014, Li et al.~2018), neural rendering
surveys, recent 3DGS papers (Mip-Splatting, 2DGS, Spec-Gaussian),
NVIDIA's CUDA--Vulkan interop sample, ANGLE\slash VKD3D as
API-translation precedents.}

% =====================================================================
\section{Conclusion}
% =====================================================================

We have argued that synthetic data generation for robotics is the wrong
workload to be running on graphics-optimized silicon, and that Vulkan is
the right API to keep. \textbf{vkSplat} is our proposed bridge: a
software Vulkan implementation whose backend is a compute-only renderer
on AI accelerators, beginning with a 3D Gaussian Splatting fast path
and growing toward a full software graphics stack. The position we
defend is that graphics APIs and graphics hardware should be decoupled,
and that the SDG market is the place where this decoupling pays off
first.

\section*{Acknowledgments}
\todo{Add acknowledgments after collaborators are named.}

\bibliographystyle{IEEEtran}
\bibliography{references}

\end{document}
